%%% 3-Phase Medical Risk Prediction Technical Roadmap (TikZ)
%%% Nature-style white background, color-blind-safe palette
%%% Compile with: xelatex fig_roadmap.tex
%%%
%%% Designed for: Chinese medical project proposals (温州市科技局/省科技厅)
%%% Archetype: architecture (3-layer vertical stack)
%%% Size: ~14.9cm wide × ~14.8cm tall (double-column compatible)
%%%
%%% Usage:
%%%   1. Copy this file, rename for your project
%%%   2. Customize phase titles, box labels, and layout coordinates
%%%   3. xelatex fig_roadmap.tex
%%%   4. pdftoppm -png -r 300 fig_roadmap.pdf fig_roadmap
%%%   5. Insert PNG into docx

\documentclass[tikz,border=1pt]{standalone}
\usepackage{xcolor,tikz}
\usepackage{fontspec}
\usepackage{xeCJK}
\setCJKmainfont{Noto Sans CJK SC}
\setCJKsansfont{Noto Sans CJK SC}
\usetikzlibrary{shapes,arrows,positioning}

% Nature 6-color palette (color-blind safe)
\definecolor{nat_blue}{HTML}{0072B2}
\definecolor{nat_cyan}{HTML}{56B4E9}
\definecolor{nat_green}{HTML}{009E73}
\definecolor{nat_orange}{HTML}{E69F00}
\definecolor{nat_pink}{HTML}{CC79A7}
\definecolor{nat_brown}{HTML}{D55E00}
\definecolor{nat_gray}{HTML}{999999}

\begin{document}
\begin{tikzpicture}[
    every node/.style={font=\sffamily},
    box/.style={draw=#1, fill=#1!10, text=#1!95!black, minimum width=22mm,
                minimum height=7mm, rounded corners=1.2mm, font=\sffamily\footnotesize,
                align=center, inner sep=2pt, line width=0.8pt},
    boxsml/.style={draw=#1, fill=#1!8, text=#1!95!black, minimum width=18mm,
                   minimum height=5mm, rounded corners=0.8mm, font=\sffamily\scriptsize,
                   align=center, inner sep=1.5pt, line width=0.5pt},
    layer/.style={draw=black!20, rounded corners=2mm, fill=black!1, line width=0.7pt},
    lbl/.style={text=black!55, font=\sffamily\small\bfseries, align=center},
    annot/.style={text=black!40, font=\sffamily\tiny, align=center},
    tag/.style={text=#1!70!black, font=\sffamily\scriptsize, align=center},
    oval/.style={draw=#1, fill=#1!10, text=#1!95!black, minimum width=20mm,
                 minimum height=6mm, rounded corners=3mm, font=\sffamily\footnotesize,
                 align=center, inner sep=2pt, line width=0.8pt},
    arr/.style={->, black!35, >=stealth, line width=0.6pt, shorten >=2pt, shorten <=2pt},
]

% ===== LAYOUT COORDINATES (modify per project) =====
% Phase 1 (指标体系构建)
\def\pOneTop{10.2}
\def\pOneBot{6.5}
\def\pOneCtr{8.35}
% Phase 2 (预测模型构建与验证)
\def\pTwoTop{5.6}
\def\pTwoBot{1.5}
\def\pTwoCtr{3.55}
% Phase 3 (护理预警与临床转化)
\def\pThreeTop{0.8}
\def\pThreeBot{-3.0}
\def\pThreeCtr{-1.1}
% Box center X positions (adjust for 3-5 elements per row)
\def\colA{-3.2}
\def\colB{0}
\def\colC{3.2}
\def\colVal{5.2}  % validation/output column
\def\colInt{1.8}  % intervention column

% ===== PHASE 1 =====
\draw[layer] (-6.5,\pOneBot) rectangle (6.5,\pOneTop);
\node[lbl] at (-5.5,\pOneTop-0.4) {第一阶段：\textcolor{nat_blue}{此处替换项目阶段1名称}};

% Flow boxes
\node[box=nat_blue] (p1a) at (\colA,\pOneCtr+0.7) {步骤A\\简要描述};
\node[box=nat_blue] (p1b) at (\colB,\pOneCtr+0.7) {步骤B\\简要描述};
\node[box=nat_blue] (p1c) at (\colC,\pOneCtr+0.7) {步骤C\\简要描述};
\draw[arr] (p1a) -- (p1b);
\draw[arr] (p1b) -- (p1c);

% Sub boxes (optional)
\node[boxsml=nat_cyan] at (\colA,\pOneCtr-0.5) {细节1};
\node[boxsml=nat_cyan] at (\colB,\pOneCtr-0.5) {细节2};
\node[boxsml=nat_cyan] at (\colC,\pOneCtr-0.5) {细节3};

% Output box (right)
\node[oval=nat_pink] (p1out) at (\colVal+0.8,\pOneCtr) {输出内容};
\draw[arr] (\colC+1.1,\pOneCtr+0.7) -- (p1out);

% ===== INTER-PHASE ARROW =====
\draw[->, nat_blue!60, >=stealth, line width=1.2pt] (0,\pOneBot+0.2) -- (0,\pOneBot-0.3);
\node[tag=nat_blue] at (0,\pOneBot-0.5) {▼ 衔接说明};

% ===== PHASE 2 =====
\draw[layer] (-6.5,\pTwoBot) rectangle (6.5,\pTwoTop);
\node[lbl] at (-5.5,\pTwoTop-0.4) {第二阶段：\textcolor{nat_green}{此处替换项目阶段2名称}};

% Left column: main flow
\node[box=nat_green] (p2a) at (\colA,\pTwoCtr+0.8) {步骤A};
\node[box=nat_green] (p2b) at (\colA,\pTwoCtr-0.1) {步骤B};
\node[box=nat_green] (p2c) at (\colA,\pTwoCtr-1.0) {步骤C};
\draw[arr] (p2a) -- (p2b);
\draw[arr] (p2b) -- (p2c);

% Right: model/analysis output
\node[box=nat_cyan] (p2d) at (\colC-0.5,\pTwoCtr-0.1) {核心输出\\建模结果};
\draw[arr] (\colA+1.1,\pTwoCtr-0.1) -- (p2d.west);

% Right: validation panel
\node[oval=nat_pink] (p2v) at (\colVal,\pTwoCtr-0.1) {验证指标\\指标A\\指标B\\指标C};
\draw[arr] (p2d.east) -- (p2v.west);

% ===== INTER-PHASE ARROW =====
\draw[->, nat_green!60, >=stealth, line width=1.2pt] (0,\pTwoBot+0.2) -- (0,\pTwoBot-0.3);
\node[tag=nat_green] at (0,\pTwoBot-0.5) {▼ 衔接说明};

% ===== PHASE 3 =====
\draw[layer] (-6.5,\pThreeBot) rectangle (6.5,\pThreeTop);
\node[lbl] at (-5.5,\pThreeTop-0.4) {第三阶段：\textcolor{nat_orange}{此处替换项目阶段3名称}};

% Main flow
\node[box=nat_orange] (p3a) at (\colA,\pThreeCtr+0.8) {分析步骤\\交互检验};
\node[box=nat_orange] (p3b) at (\colA,\pThreeCtr-0.3) {分级/分层\\标准};
\draw[arr] (p3a) -- (p3b);

% Intervention (right)
\node[box=nat_brown] (p3i1) at (\colInt,\pThreeCtr+1.0) {级别1：干预策略A};
\node[box=nat_brown] (p3i2) at (\colInt,\pThreeCtr+0.0) {级别2：干预策略B};
\node[box=nat_brown] (p3i3) at (\colInt,\pThreeCtr-1.0) {级别3：干预策略C};
\draw[arr] (p3b.east) -- (\colInt-2.0,\pThreeCtr-0.3);
\draw[arr] (\colInt,\pThreeCtr+0.5) -- (\colInt,\pThreeCtr+0.35);
\draw[arr] (\colInt,\pThreeCtr-0.4) -- (\colInt,\pThreeCtr-0.55);

% Closed loop (bottom center)
\node[oval=nat_pink] (loop) at (0,\pThreeCtr-1.6) {闭环：评估 → 分层 → 干预 → 再评估};

% ===== FIGURE TITLE =====
\node[text=black, font=\sffamily\bfseries\large] at (0,\pOneTop+0.6) {项目名称——技术路线图};
\node[text=black!50, font=\sffamily\small] at (0,\pOneTop+0.2) {关键词1 · 关键词2 · 关键词3};

% ===== FOOTER =====
\node[annot, text=black!25] at (0,\pThreeBot-0.3) {项目来源};

\end{tikzpicture}
\end{document}
